
\section{预备知识}

\subsection{Zeckendorf 数码与黄金均值语言}

取 Fibonacci 数列 $F_1=1$，$F_2=2$，$F_{k+2}=F_{k+1}+F_k$。
Zeckendorf 定理断言：每个正整数 $N$ 存在唯一展开
$$
N=\sum_{k\ge 1} c_k(N)\,F_k,\qquad c_k(N)\in\{0,1\},\quad c_k(N)c_{k+1}(N)=0.
$$
记 $\mathbf c(N)=(c_1(N),c_2(N),\dots)$ 为 Zeckendorf 数码序列。

对任意 $m\in\NN$，定义长度为 $m$ 的黄金均值（禁止词 $11$）语言
$$
\Ym=\Bigl\{y=y_1\cdots y_m\in\{0,1\}^m:\ y_i y_{i+1}=0,\ 1\le i<m\Bigr\}.
$$
经典计数给出 $\abs{\Ym}=F_{m+2}$。

\subsection{图同构、\WL{} 与相干配置（最小使用形式）}

给定两张带颜色图 $G,H$，\GI{} 判定是否存在保持颜色与邻接的顶点双射。$1$-\WL{}（颜色细化）对顶点颜色做迭代更新：每一步用“自身颜色 + 邻居颜色多重集”决定新颜色，直到稳定。稳定后的关系类可组织为（或至少诱导）相干配置，这为 Babai 框架提供了代数容器与分支控制接口 \cite{babai2016gi,ivanov2013coherent,wl2018symmetryvsregularity}。

\subsection{块系统、纤维与残余对称}

一个映射 $f:X\to Z$ 诱导分解
$$
X=\bigsqcup_{z\in Z} f^{-1}(z),
$$
其中 $f^{-1}(z)$ 称为纤维。要求同构与 $f$ 相容（例如 $f(\pi(x))=f(x)$）等价于要求置换群作用保持该分解的块结构。本文将纤维上的“不可见自由度”与 \WL{} 稳定后残余轨道对应起来，作为“时间/复杂度”的可计算对象。

